\odiseachapter{Decimoquinta parte}{Homero}

El último canto de la \emph{Odisea} es un pastiche que incluye tres pasajes diferentes.

El primero es una segunda visita al Hades, con la excusa de seguir a las almas de los pretendientes, conducidas en tropel por Hermes al prado de los asfódelos. El poema sigue insistiendo en la imagen del rebaño sacrificado, que incluso después de muertos precisa de un pastor para guiarlos al averno. No es que Homero se interese mucho por la ex-alegre pandilla de gorrones. Toda la excursión al ultramundo es una excusa para presentarnos una conversación entre Agamenón y Aquiles, en la que el primero sigue quejándose ---claramente le sentó fatal que lo asesinaran--- y de paso describe, no con poca envidia, las honras fúnebres dedicadas al segundo.

Tampoco, por supuesto, se nos cuenta nada más de las esclavas, vilmente asesinadas. ¿Es ese ninguneo una prueba de «machismo»? Al margen de las dificultades de proyectar la ética y el lenguaje moderno a una civilización separada de la nuestra por tres milenios, me parece que no es el caso. Si excluimos a Odiseo, que, todo hay que decirlo, tiene algo de \emph{Übermensch}, me atrevo a decir, y así he tratado de reflejarlo en este libro, que las auténticas protagonistas de la \emph{Odisea} son mujeres. Algunas divinas ---Atenea, Circe, Calipso---, otras humanas ---Penélope, Nausícaa, Arete, Clitemnestra--- y alguna, como Helena, en la liga especial de las semidiosas. Pero ninguna de ellas es un florero, ninguna se pliega a los deseos de los hombres, ninguna se limita a actuar como comparsa, todas se caracterizan por su firme voluntad y en la mayoría de los casos ---Penélope, Nausícaa, Arete, incluso las divinas Circe y Calipso--- por una calidad moral que ya querrían tener los hombres. En la división olímpica, aunque Zeus hace algún cameo, quien lleva absolutamente la voz cantante es Atenea. El lector simpatiza incluso con Clitemnestra, sobre todo en la versión en la que su venganza está motivada por el asesinato de su hija y resulta difícil no acabar suspirando por Helena, como aquellos bárbaros de hace tres milenios.

Así que el trato que se les da a las muchachas se entiende, precisamente, en la clave de que son esclavas. Por eso las ahorca Telémaco, por eso nada se nos dice de su viaje al ultramundo. Como tales esclavas no tienen derecho, ya no digamos a un indulto más que justo dada la levedad ---por no decir inexistencia--- de sus pecados, sino ni siquiera a una muerte «digna» ---al menos a los ojos de los que las ejecutan---. Y en ese mismo espíritu, tampoco tienen derecho a que el Poeta se fije en sus sombras.

Son esclavas, anónimas, despreciadas, ignoradas por sus amos.

Afortunadamente, en estos tiempos modernos no hay esclavitud, ni todo un continente ---África--- donde la renta per cápita anual, el acceso a la energía, al agua potable y a servicios básicos de enseñanza y salud son tan inferiores a los del primer mundo que la fortuna del gran magnate que todos conocemos rivaliza con el producto anual de las mayores economías del continente. Tampoco hay guerras innecesarias y crueles, ni genocidios, ni desigualdades entre los adinerados y los desposeídos, ni millones de personas dispuestas a jugarse la vida por el espejismo de una vida mejor, es decir, esclavos que se suben a la patera como las muchachas de la casa de Ulises se subían al lecho de los pretendientes, porque no tienen nada mejor a lo que agarrarse.

Quizás por eso nos escandalizamos tanto frente al abuso. Tanto, que ninguna película lo muestra, porque barrer la ignominia bajo la alfombra es algo a lo que estamos muy bien acostumbrados.

El segundo fragmento del canto XXIV es la visita a Laertes. A pesar de que se trata de un encuentro padre-hijo y a pesar de mi propia referencia emocional ---he tenido la fortuna de tener unos padres y unos hijos maravillosos---, el episodio no consiguió emocionarme en mi primera lectura, ni tampoco me ha emocionado en las lecturas posteriores. Quizás un poema reciente me permita explicarme un poco. Se llama «Lecciones de conducir»:

\begin{verse}
A papá le preocupaba que aprendiera a conducir,\\
le preocupaba tanto, que tomamos la costumbre,\\
de hacer prácticas, mucho antes,\\
de tener la edad para sacarme el carnet.

Aquellas escapadas,\\
maravillosamente ilegales, mi yo adolescente,\\
al timón de nuestro enorme coche familiar,\\
más grande que el \emph{Nautilus}.

Era también nuestro pequeño secreto,\\
nos fugábamos discretamente,\\
por carreteras soleadas que discurrían\\
entre suburbios industriales y aldeas escondidas.

Parábamos en pequeños cafés,\\
donde charlábamos largamente, después de todo,\\
éramos socios al volante y lo suyo,\\
era intercambiar historias.

Y mi padre tenía,\\
tantas historias que contarme.

Cuando saqué el carnet siguió insistiendo,\\
en que practicar era imprescindible,\\
declarándose todavía francamente inquieto,\\
con mi obvia falta de reflejos.

Era una extraña obsesión, de hecho,\\
yo había sacado el carnet a la primera,\\
y después de tanta práctica, sabía muy bien\\
que me sobraba habilidad. No importaba,

papá declaró que todavía eran necesarias,\\
veinte mil leguas de viaje.

Así que las hicimos. Viajábamos a veces,\\
al pueblo veraniego, pernoctábamos,\\
en la casita que teníamos allí,\\
antes de volver a casa al día siguiente.

Papá cancelaba sus clases,\\
anulaba compromisos, se saltaba citas,\\
todo fuera por enseñarme,\\
a conducir, decía, como Dios manda.

Y mamá, de repente,\\
parecía igualmente preocupada,\\
y declaraba que aquellas excursiones,\\
eran absolutamente necesarias.

Ayer, pasé el día junto a mi hijo,\\
en otro largo viaje, conduce tan bien\\
como solía hacerlo yo, y aun así,\\
prolongaría ese viaje hasta Finisterre.

Viajaría con él veinte mil leguas, hasta llegar,\\
al final del Arco Iris, donde papá nos espera.
\end{verse}

Mi padre no era de los que te abrazan cada vez que pueden ---yo, en cambio, sí y mis pobres hijos tienen que soportar al pesado de su viejo achuchándoles con cualquier excusa--- ni era «mi mejor amigo», ni, que yo sepa, me dijo nunca «te quiero» con todas las letras ---también en eso somos muy diferentes, yo se lo repito a mis hijos hasta agotarlos---. Pero el poema no exagera en lo que cuenta y es solo una mínima parte de lo que hacía por mí y por mis hermanos y hermanas.

Quizás sea un prejuicio mío ---el sesgo emocional que produce darte cuenta de que tu padre ha cancelado una reunión importante para darte una clase de conducir que no necesitas--- lo que me deja frío frente a unos versos formalmente magníficos pero tan formales que uno no acaba de ver el sentimiento:

\begin{verse}
«Aquel por el que preguntas soy yo, padre mío, yo mismo,\\
que he regresado a la patria, después de veinte años ido.\\
Así que ten ya tu llanto, ten tú el lagrimoso gemido,\\
que ahora te tengo que hablar, y afanarse en el trance es preciso:\\
en nuestra casa he matado a los pretendientes altivos\\
dando a su ultraje doliente y a todos sus males castigo».
\end{verse}

\begin{verse}
Y le decía Laertes, alzando la voz, en respuesta:\\
«Si de verdad tú Odiseo, mi hijo tú, que aquí llegas,\\
dame ya seña bien clara para que yo me aconvenza».
\end{verse}

\begin{verse}
Y respondiendo decíale así el milardido Odiseo:\\
Pon esos ojos aquí, a esta herida no más lo primero,\\
la que me abrió con su blanco colmillo en Parnaso aquel puerco\\
cuando allí fui porque tú y la alta madre al abuelo materno,\\
a Autólico, me encomendasteis en pos de los dones aquellos\\
que prometió cuando vino y me confirmó con su gesto.
\end{verse}

Laertes exige una prueba y Ulises le muestra la cicatriz que dejó el jabalí y que también ha reconocido antes su niñera ---viene a ser, en el poema, como su DNI digital---. También le habla de los árboles del huerto, un pasaje harto más íntimo que el de la cansina cicatriz:

\begin{verse}
Y, sí, también te diré de los árboles de este buen huerto,\\
los que me diste una vez, que uno a uno los iba pidiendo\\
yo cuando niño, y pegado a ti andaba detrás entre ellos\\
por el hortal; tú decías sus nombres y yo iba aprendiendo.\\
Trece perales me diste, y diez manzanos con ellos,\\
de higueras cuatro decenas, y liños en el viñedo\\
tú me darías cincuenta, en sazón cada cual a su tiempo,\\
que aquí racimos de toda sazón se levantan del suelo\\
cuando de arriba los cargan las estaciones de Zeus.
\end{verse}

\begin{verse}
Tal dijo; y ya sus rodillas y su corazón flaqueaban,\\
reconociendo señal tras señal que Odiseo le daba.\\
Echó a su hijo los brazos, y contra su cuerpo apretaba\\
el sufrimundo Odiseo a su padre, que allí desmayaba.
\end{verse}

Bonito, pero Homero, como mi padre, no es muy amigo de expansiones paterno-filiales y enseguida Laertes pasa a asuntos prácticos, preocupándose por la posible revuelta de los familiares de los pretendientes masacrados. A partir de ahí se da por zanjada la parte sentimental y vuelven a aparecer los guerreros. En el tercer fragmento, Laertes toma parte en el último enfrentamiento entre Ulises y sus fieles y la banda de enojados padres a los que acaban de matarles a sus hijos. Y también en este caso me he encontrado yo siempre del lado equivocado. Si hubiera sido mi hijo el asesinado me habría faltado tiempo para tomar una lanza y enfrentarme al asesino. Sin duda mi destino habría sido el mismo que el de Eupites, el padre de Antínoo, que se pone a la cabeza de la partida de rebeldes:

\begin{verse}
Eupites los conducía igual que a niños pequeños;\\
y él, que la muerte del hijo pensaba vengar, nunca luego\\
ya tornaría, que allí lo aguardaba el destino funesto.
\end{verse}

Me habría pasado como a Eupites, que muere al enfrentarse a Laertes, quien, por cierto, debe haber fumado algo muy fuerte: ha pasado de ser un viejo deshecho a estar en plena forma y sentir su pecho henchido de orgullo:

\begin{verse}
[...] y en gozo Laertes a sí se decía:\\
«Dioses queridos, ¡qué día me dais! ¡Y cuánta alegría\\
ver a mi hijo y su hijo por la bravura en porfía!».
\end{verse}

En realidad, no necesita fumar nada, tiene un camello divino para darle un chute de olímpica energía:

\begin{verse}
Tal dijo; y Palas Atena alentole un vigor sin medida.\\
Y, haciendo rezo en voz alta a la de Zeus Cronida,\\
lanzó, blandiendo con pulso, su azcona larguisombría\\
y al casco le traspasó de Eupites la broncicarrilla\\
que de asta no lo libró, y entró el rejo de claro en su crisma:\\
él al caer retumbó, y sus armas chocáronle encima.
\end{verse}

¡Pobre Eupites! Odiseo y los suyos caen sobre el resto de la banda, con la sana intención de cortarlos en cubitos, pero los olímpicos ya se han divertido bastante y decretan el fin de las hostilidades. Atenea da la voz de basta:

\begin{verse}
«Teneos, hombres de Ítaca, de la agria guerra, y sin sangre\\
marchad cada uno a lo vuestro concordes y cuanto antes».
\end{verse}

\begin{verse}
Dijo Atenea, y a aquellos tomó la pavura amarilla;\\
y de sus manos, medrosos, volaban las armas, que iban\\
todas al suelo a parar, a la voz de la diosa que oían;\\
y se volvían a la ciudad anhelosos de vida.\\
El sufrimundo Odiseo rugió en malhorrísona grita\\
y prieto en sí se fue a ellos como águila de altanería.\\
Entonces suelta el Cronión una humosa centella encendida,\\
que vino a dar a los pies de la patrisoberbia ojiviva.
\end{verse}

\begin{verse}
Y dijo Atena ojiglauca, volviendo a Odiseo su vista:\\
«Sangre Laercia, Odiseo milmañas, hijo del cielo,\\
tente y pon fin al azote común de la guerra y su duelo,\\
para que no se te enoje Zeus Cronión miraluengo».
\end{verse}

Así que la \emph{Odisea} acaba, literalmente, con un deus ex machina. Mi bienquerido director, dicho sea de paso, debería haberse estudiado el final del canto XXIV, con lo que nos habría ahorrado las ridículas frases grandilocuentes que pronuncian al alimón Penélope y Ulises mientras navegan hacia el sol poniente. Homero, en cambio, sabe que lo bueno, si breve. Después de la intervención de Atenea, cierra de un plumazo:

\begin{verse}
Tal dijo Atena, y él bien la escuchó, en el alma alegrándose.\\
Y luego entre ellos dispuso los juros de paz perdurable\\
la diosa Palas Atena, la hija de Zeus cabrillante,\\
a Méntor siendo en figura y a par en voz semejante.
\end{verse}

Y así acaba la \emph{Odisea}, la obra magna de Homero.

Si es que hubo un Homero, claro. Infinitas son las disputas de los magos y más infinitos aún los debates sobre si Homero fue una sola persona o un colectivo que edita y unifica una rica tradición oral anterior\footnote{Una introducción a la cuestión homérica puede encontrarse en Gregory Nagy, \emph{Homeric Questions}: \url{https://chs.harvard.edu/book/nagy-gregory-homeric-questions/}}.
%También se discute vivamente si la Ilíada y a Odisea pertenecen al mismo autor.

La noción de que ambos poemas épicos se derivan de una amplia ---y muy antigua--- tradición oral tiene ya cien años de antigüedad. El personaje más importante asociado a la idea es Milman Parry (1902--1935), quien ---nobleza obliga--- tuvo un final prematuro y absurdo, al dispararse por accidente su propio revólver. Imposible no pensar en los caprichosos dioses.

La idea de Parry resulta muy convincente para cualquiera que haya visto en acción a los bertsolaris de Euskadi, o se haya deleitado con las «peleas de gallos», en las que jóvenes raperos se dan bofetadas verbales, inventando sus versos al vuelo\footnote{Mi hijo es muy aficionado a estas peleas de gallos y hace poco escuchábamos una en la que se esgrimían con descaro estas líneas: «Tú usas la violencia / yo soy inteligente / ¡y te hago bullying... mentalmente!»}. Unos y otros improvisan, sí, pero basándose en patrones bien establecidos y memorizados que les permiten ejecutar en tiempo real variaciones sobre un tema y les facilitan el ritmo y la rima. Es algo parecido a lo que se hace en jazz, donde la inagotable creatividad de los músicos se ancla en un tronco musical muy robusto.

Parry argumentó exactamente eso, mostrando que tanto la \emph{Ilíada} como la \emph{Odisea} usan fórmulas orales, temas fijos y escenas más o menos sincopadas cuyo objetivo es permitir la improvisación al vuelo. A lo largo de estas páginas hemos visto muchas de estas fórmulas. Una de mis favoritas se usa para nombrar a Ulises:

\begin{verse}
Sangre Laercia, Odiseo milmañas, hijo del cielo
\end{verse}

Ya vimos que las fórmulas de cortesía son idénticas. Méntor usa la misma para agasajar a Telémaco que Polifemo para burlarse de los aqueos. Los ejemplos abundan. Atenea es ojiglauca, Zeus es cabrillante, Helena es mantodevuelo y cabellilinda, Circe beltrenza y vocihumana, Penélope milpensamientos. Todo un armazón de trucos nemotécnicos y fórmulas prefabricadas disponible para que el bardo pueda memorizar e improvisar. De hecho, Parry llamó a estos patrones fijos, precisamente, «fórmulas». Ya hemos visto algunas, ahí van algunas más: \gr{ῥοδοδάκτυλος Ἠώς} (rhododáktylos Ēṓs) o «Aurora de rosados dedos»; \gr{ἔπεα πτερόεντα} (épea pteróenta) o «aladas palabras», \gr{πολύμητις Ὀδυσσεύς} (polýmētis Odysseús) que Tobal traduce como Odiseo mentisagaz o milastucias, \gr{ἕρκος ὀδόντων} (hérkos odóntōn) o «el cerco de los dientes». Estas fórmulas, adaptadas además a la métrica de los hexámetros dactílicos, permitían componer al vuelo:

\begin{verse}
\gr{ἦμος δ᾽ ἠριγένεια φάνη ῥοδοδάκτυλος Ἠώς}\\
êmos d' ērigéneia phánē rhododáktylos Ēṓs\\
«Cuando se mostró la que nace temprano, Aurora de rosados dedos»
\end{verse}

\begin{verse}
\gr{καί μιν φωνήσασ᾽ ἔπεα πτερόεντα προσηύδα}\\
kaí min phōnḗsas' épea pteróenta prosēúda\\
«Y alzando la voz le dirigió aladas palabras»
\end{verse}

\begin{verse}
\gr{τὴν δ᾽ ἀπαμειβόμενος προσέφη πολύμητις Ὀδυσσεύς}\\
tḕn d' apameibómenos proséphē polýmētis Odysseús\\
«Respondiéndole, dijo el muy astuto Odiseo»
\end{verse}

\begin{verse}
\gr{τέκνον ἐμόν, ποῖόν σε ἔπος φύγεν ἕρκος ὀδόντων}\\
téknon emón, poîón se épos phýgen hérkos odóntōn\\
«Hija mía, ¿qué palabra se te escapó del cerco de los dientes?»
\end{verse}

La fértil intuición de Parry ha sido extendida y cuantificada a lo largo de las décadas, incluyendo estudios estadísticos que muestran que la densidad de «fórmulas» es muy superior a la que se encuentra en autores clásicos posteriores, como Eurípides o Virgilio.

¿Qué demuestra entonces el análisis de Parry? Por un lado, en el debate anterior a su trabajo, un argumento repetido a favor del concepto de los múltiples autores era el aparente batiburrillo de estilos y capas dialectales, que sugería que el poema había sido compuesto por varios bardos desparramados en varias épocas. Milman metisagaz demuestra que repeticiones, dobladillos, cantinelas y remiendos eran en realidad parte de la técnica de composición. Homero podría seguir siendo un solo autor, posiblemente ---pero no necesariamente--- literato, que recoge la rica tradición heredada en forma de relatos fragmentados y le da coherencia, creando una épica unitaria.

Homero sería entonces la voz en cuyo eco resuenan muchas voces, un autor único y a la vez plural. Puede no ser la única interpretación válida, pero debo confesar aquí que es la que más me gusta. Resuena con esa idea romántica a la que tantos nos asimos como el divino Odiseo a los maderos tras el naufragio. \emph{La vida es corta, pero el arte es inmortal}.
